\section{Conclusion}

\begin{frame}{Summary of this Work}{}

    \textbf{IFAC}

    \begin{itemize}
        \item Currently, effective-space formulation and input constraints are considered.
        \item The input constraint will be removed and the effective-space formulation will be more rigorously justified.
    \end{itemize}

    \textbf{Future work}

    \begin{itemize}
        \item Additional constraints, such as state constraints and CBF-motivated constraints, and input constraints will be considered.
        \item Real-time implementation of the proposed method will be investigated.
    \end{itemize}

    \textbf{Towards Deep Learning}

    \begin{itemize}
        \item Using obtained $\mv{y}$, we can extract the contraction metric $\mm{M}$.
        \item Using neural networks, we can learn the mapping from the state $\mv{x}$ to the contraction metric $\mm{M}$.
        \item Hence, if we can justify that the obtained metric is a valid contraction metric under imposed constraints, we can use the learned metric for control design and analysis.
    \end{itemize}
    
\end{frame}

